%! TeX program = lualatex
\usepackage{luatexja}
\usepackage[match]{luatexja-fontspec}
%\jfont{noto-jp}
\usepackage[noto]{luatexja-preset}
\usepackage[x11names,dvipsnames,table]{xcolor}
% \usepackage[utf8]{inputenc}
% \usepackage[]{url}
\usepackage{graphicx}
\graphicspath{ {./figures/} }
\usepackage[]{float}
\usepackage[]{verbatimbox}
\usepackage[listings,breakable,skins,theorems]{tcolorbox}
\usepackage[]{listings}
%\usepackage[style=numeric]{biblatex}


% select font family
\renewcommand*\familydefault{\sfdefault}
\usepackage{textcomp}

\usepackage{booktabs}
\usepackage{enumitem}
% Hide page number on empty pages
\usepackage{emptypage}
\usepackage[marginparwidth=75pt]{geometry}

\usepackage{subcaption} % arbitrary captions
% tabular packages
\usepackage{multicol} % merged columns in tabular environments
\usepackage[]{multirow}
\usepackage[]{tablefootnote}

% tikz packages
\usepackage[]{tikz}
\usetikzlibrary{calc,matrix,shapes,positioning}

\usepackage[normalem]{ulem}

% Extraneous symbols
\usepackage[misc]{ifsym}
\usepackage{pifont} % access to dingbats style symbols

% Maths related packages
\usepackage[]{cancel} % allows for writing 
\usepackage{amsmath, amsfonts, mathtools, amsthm, amssymb}
\usepackage[]{derivative}

% Script capitals 
\usepackage{mathrsfs}
% Bold math
\usepackage{bm}
%Some useful shortcuts for math environment Fonts
\newcommand\N{\ensuremath{\mathbb{N}}}
\newcommand\R{\ensuremath{\mathbb{R}}}
\newcommand\Z{\ensuremath{\mathbb{Z}}}
\renewcommand\O{\ensuremath{\emptyset}}
\newcommand\Q{\ensuremath{\mathbb{Q}}}
\newcommand\C{\ensuremath{\mathbb{C}}}
\newcommand\Px{\ensuremath{\mathbb{P}}}

%%%%%%%%%%%%%%%%%%%%%%%%%%%%
%%% USER DEFINED COLOURS %%%
%%%%%%%%%%%%%%%%%%%%%%%%%%%%
\colorlet{b}{blue}
\colorlet{B}{RoyalBlue}
\colorlet{g}{green!50!black}
\colorlet{r}{red!100}
\colorlet{p}{purple}
\colorlet{m}{Magenta}
\colorlet{o}{orange}
\colorlet{fb}{Firebrick3}
\definecolor{t1}{HTML}{E53935}
\definecolor{t2}{HTML}{C8A200}
\definecolor{t3}{HTML}{43A047}
\definecolor{t4}{HTML}{1E88E5}
\definecolor{t5}{HTML}{8e24aa}
\definecolor{t6}{HTML}{78909C}
%\definecolor{tn}{HTML}{red!100}


%%%%%%%%%%%%%%%%%%%%%%%%%%%%%
%%% USER DEFINED COMMANDS %%%
%%%%%%%%%%%%%%%%%%%%%%%%%%%%%

\newcommand{\norm}[1]{\left \lVert #1 \right \rVert}
\newcommand{\arrowheadbullet}{\ding{228} }

% Enum item custom series
\newcommand*{\encircled}[1]{%
  \footnotesize\protect\tikz[baseline=-3px]%
  \protect\node[fill=gray!50, shape=circle, draw=none,%
inner sep=1.2pt](n1){#1};
}
% Enum item custom series : Rounded Square with Capital Letters
\newcommand*{\ensquared}[1]{%
  \large\protect\tikz[baseline=-3px]\normalsize%
  %\footnotesize\protect\tikz[baseline=-3px]%
  \protect\node[fill=none, shape= rectangle, draw=black, rounded corners=0.1cm,%
  inner sep=6pt](n1){\textbf{#1}};
}
% Environments
\makeatother

% framed envs
% \usepackage{thmtools}
%\usepackage[framemethod=TikZ]{mdframed}
% \mdfsetup{skipabove=1em,skipbelow=0em}

\pgfdeclarehorizontalshading{defbackground}{100bp}
  {color(0bp)=(white);color(100bp)=(ForestGreen!45)}

\pgfdeclarehorizontalshading{propbackground}{100bp}
  {color(0bp)=(white);color(100bp)=(NavyBlue!45)}



%-----------------------------
% Listing Styles for Languages
%-----------------------------
% Python
\definecolor{comment}{rgb}{0.58,0,0.82}
\definecolor{linenumber}{rgb}{0.5,0.5,0.5}
\definecolor{string}{RGB}{229,88,56}
\lstdefinestyle{py3}{
%  backgroundcolor=\color{ForestGreen!15!white},
  commentstyle=\color{comment},
  keywordstyle=\color{magenta},
  numberstyle=\tiny\color{linenumber},
  stringstyle=\color{string},
  basicstyle=\ttfamily\footnotesize,
  breakatwhitespace=false,
  breaklines=true,
  captionpos=b,
  keepspaces=true,
  morecomment=[is]{\#\%}{\#:}
  % numbers=left,
  numbersep=5pt,
  showspaces=false,
  showstringspaces=false,
  showtabs=true,
  tabsize=4,
  language=Python
}
\lstdefinestyle{general}{
  backgroundcolor=\color{ForestGreen!15!white},
  commentstyle=\color{comment},
  keywordstyle=\color{magenta},
  numberstyle=\tiny\color{linenumber},
  stringstyle=\color{string},
  basicstyle=\ttfamily\footnotesize,
  breakatwhitespace=false,
  breaklines=true,
  captionpos=b,
  keepspaces=true,
  % numbers=left,
  numbersep=5pt,
  showspaces=false,
  showstringspaces=false,
  showtabs=false,
  tabsize=2
}
% End environments 

% Hide year part of lecture command
% hide parts
  \newcommand\hide[1]{}

  \makeatletter
  \usepackage{xifthen}
  \def\testdateparts#1{\dateparts#1\relax}
  \def\dateparts#1 #2 #3 #4 #5\relax{
    \marginpar{\small\textsf{\mbox{#1 #2 #3 #5}}}
  }
  \def\@lecture{}%
  \newcommand{\lecture}[3]{%
    \ifthenelse{\isempty{#3}}{%
      \def\@lecture{Lecture #1}%
    } {%
      \def\@lecture{Lecture #1: #3}%
    }%
    \subsection*{\@lecture}
    \testdateparts{#2}
   % \marginpar[\scriptsize\textsf{\mbox{#2}}]{\scriptsize\textsf{\mbox{#2}}}
  }
  \def\@tutorial{}%
  \newcommand{\tutorial}[3]{%
    \ifthenelse{\isempty{#3}}{%
      \def\@tutorial{Tutorial #1}%
    } {%
      \def\@tutorial{Tutorial #1: #3}%
    }%
    \subsubsection*{\@tutorial}
    \testdateparts{#2}
   % \marginpar{\small\textsf{\mbox{#2}}}
  }
%  \def\@seminar{}%
%  \newcommand{\seminar}[3]{%
%    \ifthenelse{\isempty{#3}}{%
%      \def\@seminar{Seminar #1}%
%    } {%
%      \def\@seminar{Seminar #1: #3}%
%    }%
%    \subsection*{\@seminar}
%    \testdateparts{#2}
%   % \marginpar{\small\textsf{\mbox{#2}}}
%  }
  \renewcommand{\chaptername}{Lesson} % customize chapter names
  \def\@topic{}%
  \newcommand{\topic}[1]{%
    \def\@topic{#1}%

    \chapter{\@topic}
  }



% \def\@date{}%
%  \newcommand{\wdate}[3]{%
%    \ifthenelse{\isempty{#3}}{%
%      \def\@date{Date #1}%
%    } {%
%      \def\@date{Date #1: #3}%
%    }%
%    \subsection*{\@date}
%    % \marginpar{\small\textsf{\mbox{#2}}}
%    \testdateparts{#2}
%  }
% 
% counter for exercises in a lecture
  \newcounter{exercise}[chapter] % resets with each new chapter in document 
  \renewcommand{\theexercise}{\thechapter.\arabic{exercise}}
% % Other commands to set
% 
% \newcommand{\subexercise}[1]{%
%   \subsubsection*{Exercise \@exercise.#1}
% }
% \def\@tutorial{}%
% \newcommand{\tutorial}[3]{
%   \ifthenelse{\isempty{#3}}{%
%     \def\@tutorial{Tutorial #1}%
%   }{%
%     \def\@tutorial{Tutorial #1: #3}%
%   }%
%   \subsection*{\@tutorial}
%   \testdateparts{#2}
% }
% 
% \renewcommand{\theequation}{\thechapter.\arabic{equation}}
% 
% New command for double underlined text in math env
  \newcommand{\ulinedbl}[1]{%
    \underline{\underline{\textrm{ #1 }}} 
  }


  \renewcommand{\theequation}{\thechapter.\arabic{equation}}

% Homework insertion
  \pgfdeclarehorizontalshading{hwbackground}{100bp}
  {color(0bp)=(RoyalPurple!25);color(100bp)=(black!5)}
  %\newmdenv[
  %linecolor=Gray!70!Black, linewidth=1pt,%
  %roundcorner=3pt,%
  %everyline=true,
  %splittopskip=20pt,
  %apptotikzsetting={\tikzset{mdfbackground/.append style ={%
  %    shading=hwbackground},
  %    mdfframetitlerule/.append style={%
  %    linewidth=0pt},
  %    mdfframetitlebackground/.append style={fill=none}
  %}},
  %frametitlerule=false
  %]{homework}
  %\newcommand{\hw}[2]{
  %  \begin{homework}[frametitle=#1]
  %    \ding{228} #2
  %  \end{homework}
  %}
% Create color shorthand for background color
\colorlet{codebg}{ForestGreen!15}
% Further shadings

% Tcolorbox alternative to mdenv

\newcounter{homework}[chapter] % resets with each new chapter in document 
\NewTColorBox[use counter=homework]{homework}{ O{}}{
  enhanced jigsaw,
  %frame hidden,
  colframe = Gray!70!Black,
  coltitle = RoyalPurple!5!white, fonttitle = \bfseries,
  description font = \mdseries,
  title = #1,
  colframe=Purple4!75!black,
  interior style = {left color={RoyalPurple!20!white}, right color={RoyalPurple!10}}
}
  \newtcolorbox[use counter=exercise, number within=chapter]{exbox}[1]{
    enhanced jigsaw,
    fonttitle=\Large\bfseries,coltitle=black,
    colframe=RoyalPurple!15!white,title=Ex. ~\thetcbcounter: #1, breakable=true,
    opacityback=0,opacitybacktitle=0,    
    extras={frame empty, interior empty},
  }
  \NewTCBListing[auto counter, number within=chapter]{code}{O{}O{}}{
    enhanced jigsaw,
    breakable=true,
    frame hidden,
    title style={left color=RoyalPurple!25, right color={codebg!0}},
    titlerule=0mm,
    coltitle=black,
    interior style={left color ={codebg},right color = {codebg}},
    borderline={1.5pt}{.1pt}{white},
    listing options={language=Python,style=py3,numbers=none},
    listing only,
    title=\textbf{Listing~\thetcbcounter},
    after title={: #1},
    #2
  }
  % 3 optional args: [general options][background color]{input listing file}[title] 
  \NewTCBInputListing[use counter from=code, number within=chapter, list inside=code]{\codeInput}{O{}O{codebg} m O{} }{
    enhanced jigsaw,
    breakable=true,
    frame hidden,
    title style={left color=RoyalPurple!25, right color={#2!0}},
    titlerule=0mm,
    interior style={left color ={#2},right color = {#2}},
    borderline={1.30pt}{.1pt}{white},
    borderline east={1.5pt}{.1pt}{white},
    listing file={#3},
    listing only,
    title=\textbf{Listing~\thetcbcounter},
    coltitle=black,
    after title={: #4},
    #1  
  }
  % Replacement for mdenvtheorem environments using tcolorbox
  % \NewTcbTheorem[⟨init options⟩]{⟨name⟩}{⟨display name⟩}{⟨options⟩}{⟨prefix⟩}
  \NewTcbTheorem[auto counter,list inside=definitions]{definition}{Definition}{
    enhanced jigsaw,
    frame hidden,
    theorem style=break,
    rightrule=0mm, 
    sharp corners,
    leftrule=1pt,
    toprule=0pt,
    bottomrule=0pt,
    coltitle = ForestGreen!100!black!75, fonttitle = \bfseries,
    description font = \mdseries,
    interior style={left color={ForestGreen!10!white}, right color={ForestGreen!30}},
    borderline west={2pt}{0pt}{ForestGreen!70} 
  }{def}
  % New property theorem
  \NewTcbTheorem[auto counter, list inside=properties]{property}{Property}{
    enhanced jigsaw,
    frame hidden,
    theorem style=break,
    rightrule=0mm, 
    sharp corners,
    leftrule=1pt,
    toprule=0pt,
    bottomrule=0pt,
    leftrule=2pt,
    coltitle = NavyBlue,
    fonttitle = \bfseries,
    description font = \mdseries,
    interior style={left color={NavyBlue!10!white}, right color={NavyBlue!30}},
    borderline west={2pt}{0pt}{NavyBlue!70} 
  }{prop}
  % New theorem env
  \NewTcbTheorem[auto counter, list inside=theorems]{theorem}{Theorem}{
    enhanced jigsaw,
    frame hidden,
    theorem style=break,
    rightrule=0mm, 
    sharp corners,
    leftrule=1pt,
    toprule=0pt,
    bottomrule=0pt,
    coltitle = RawSienna,
    fonttitle = \bfseries,
    description font = \mdseries,
    interior style={left color={RawSienna!10!white}, right color={RawSienna!30}},
    borderline west={2pt}{0pt}{RawSienna!70} 
  }{th}
  \NewTcbTheorem[auto counter, list inside=grammar]{grammar}{Grammar}{
    enhanced jigsaw,
    frame hidden,
    theorem style=break,
    rightrule=0mm, 
    sharp corners,
    leftrule=1pt,
    toprule=0pt,
    bottomrule=0pt,
    coltitle = orange,
    fonttitle = \bfseries,
    description font = \mdseries,
    interior style={left color={orange!10!white}, right color={orange!30}},
    borderline west={2pt}{0pt}{orange!70} 
  }{grm} 
  \NewTcbTheorem[auto counter, list inside=notes]{note}{Note}{
    enhanced jigsaw,
    frame hidden,
    theorem style=break,
    rightrule=0mm, 
    sharp corners,
    leftrule=1pt,
    toprule=0pt,
    bottomrule=0pt,
    coltitle = Honeydew4,
    fonttitle = \bfseries,
    description font = \mdseries,
    interior style={left color={Honeydew3!10!white}, right color={Honeydew3!30}},
    borderline west={2pt}{0pt}{Honeydew3!70} 
  }{not}
  \NewTcbTheorem[auto counter, list inside=vocab]{vocab}{Vocab}{
    enhanced jigsaw,
    frame hidden,
    theorem style=break,
    rightrule=0mm, 
    sharp corners,
    leftrule=1pt,
    toprule=0pt,
    bottomrule=0pt,
    coltitle = DarkGoldenrod2,
    fonttitle = \bfseries,
    description font = \mdseries,
    interior style={left color={DarkGoldenrod2!10!white}, right color={DarkGoldenrod2!30}},
    borderline west={2pt}{0pt}{DarkGoldenrod2!70} 
  }{voc}% Exercise env
  \newcommand{\mdex}[2]{
    \begin{exbox}{#1}
      #2
    \end{exbox}
  }
  \def\@exercise{}%
  \newcommand{\exercise}[2]{%
    \refstepcounter{exercise} % increment exercise number before printing it
    \ifthenelse{\isempty{#2}}{
      \def\@exercise{Exercise #1}%
    }  {%
      \def\@exercise{Exercise #1: #2}%  
    }
    \subsubsection*{\@exercise}
    % \testdateparts{#2}
  }

  % Script dialogue formatting
  \newenvironment{dialogue}[2][black] {%
    \begin{center}
      \textcolor{#1}{
      \textbf{\Large#2}}\\\vskip1em
      \begin{minipage}[b]{0.55\textwidth}
        \raggedright
      }{%
      \end{minipage}
    \end{center}
  }
  % Fancy headers and footers
  \usepackage{fancyhdr}
  \pagestyle{fancy}
  \fancyhead{} % clears all pre-existing fields
  % name for open-book exam notes
  % \fancyhead[LE,RO]{Hugh Lam}
  % for twoside styling
  % \fancyhead[RO,LE]{\@lecture}  % Right odd, Left even
  % 
  % \fancyfoot[RO,LE]{\thepage}   % Right odd, Left even
  % \fancyfoot[RE,LO]{}           % Right even, Left odd
  % \fancyfoot[C]{\leftmark}      % Center

  \fancyhead[R]{\@lecture}  % Right odd, Left even
  \fancyhead[L]{\thesection}
  \fancyfoot[R]{\thepage}   % Right odd, Left even
  \fancyfoot[C]{\leftmark}      % Center


  \makeatother

  % Figure support as explained in my blog pos
  \usepackage{import}
  \usepackage{xifthen}
  \usepackage{pdfpages}
  \usepackage{transparent}
  \newcommand{\incfig}[1]{%
    \def\svgwidth{\columnwidth}
    \import{./figures/}{#1.pdf_tex}
  }
  \newcommand{\incvarfig}[2][\columnwidth]{
    \def\svgwidth{#1}
    \import{./figures/}{#2.pdf_tex}
  }
  
  % alternative command for .eps files
  
  \newcommand{\inceps}[2][\columnwidth]{
    \def\svgwidth{#1}
    \import{./figures/}{#2.eps_tex}
  }
  
  % Fix some stuff
  % %http://tex.stackexchange.com/questions/76273/multiple-pdfs-with-page-group-included-in-a-single-page-warning

  \author{Hugh Lam}

  % referencing and linking capabilities (must be last package loaded)

  \usepackage{hyperref}
  %setup hyperref formatting
  \hypersetup{
    colorlinks=true,
    linkcolor=DarkOrchid1,
    filecolor=RoyalBlue,
    urlcolor=Magenta
  }% Some basic packages

